\documentclass[10pt]{article}

\usepackage[utf8]{inputenc}

\usepackage[T1]{fontenc}

\usepackage{amsmath}

\usepackage{amsfonts}

\usepackage{amssymb}

\usepackage[version=4]{mhchem}

\usepackage{stmaryrd}

\usepackage{graphicx}

\usepackage[export]{adjustbox}

\graphicspath{ {./images/} }



\begin{document}

с запоминанием их текущих состояний и с последующей активацией или возобновлением др. процессов.



Наиболее известными и простыми программными механизмами синхронизации являются семафоры и события. С е м а ф о р - это специальная управляющая переменная, принимающая целочисленные значения. Семафор обычно связан с нек-рым конфликтным ресурсом. $К$ семафору применимы только две операции $P$ и $V$. Если в ходе исполнения процесса встретится операция $P(s)$, где $s$-семафор, то процесс может продолжаться с уменьшением значения $s$ на 1 только в том случае, когда значение $s$ положительно; в противном случае он приостанавливается и занимает место в очереди $q(s)$ процессов, ждущих соответствующего ресурса. Операция $V$ увеличивает значение $s$ на 1 и возобновляет первый в очереди $q(s)$ процесс. Механизм семафоров широко используется в языках управления процессами в операционных системах ӘВМ и в ряде универсальных языков программирования (напр., алгол68). Механизм с о б ы т и й включает управляющие переменные, текущие значения к-рых отмечают наступление каких-либо программных или системных событий (завершение процессов, прерывания и т. п.), и специальные операторы ожидания событий



Семафоры и события являются универсальными средствами синхронизации, но они слишком примитивны, и неправильное их использование может привести к аварийным ситуациям таким, как взаимная блокировка процессов (напр., два процесса требуют для своего продолжения по два ресурса и каждый «захватил» по одному). Стремление повысить надежность программирования приводит к появлению более сложных механизмов синхронизации: «п о ч т о в ы е я щ и к и»особые структуры для обмена сообщениями, для к-рых фиксированы правила работы с ними параллельных процессов; м о н и т о р - наборы процедур и данных, к к-рым процессы могут обрашаться только поочередно и к-рые содержет заданные программистом правила организации взаимодействий.



Чисто асинх ронное П. П. используется для организации вычислений в распределенных вычислительных спстемах, в к-рых полностью исключены конфорикты по ресурсам Стремление упростить организацию взаимодействий между процессами с общими ресурсами привлекает внимание к асинхронным методам вычислений, в к-рых разрешен нерегламентированный доступ параллельных процессов к общим ресурсам. Напр., разрабатываются асинхронные алгоритмы, в к-рых параллельные процессы обмениваются данными с общей памятью, причем неупорядоченный доступ к памяти не мешает достижению однозначного результата.



В теории П. п. разрабатываются формальные модели параллельных программ, процессов и систем и с их помощью исследуются различные аспекты П. п.: автоматич. распараллеливание последовательных алгоритмов и программ; разработка параллельных методов вычислений для разных классов задач и разных классов параллельных вычислительных систем; оптимальное планирование параллельных вычислений и распределение ресурсов между параллельными процессами; формальное описание семантики программ и языков П. п. Среди таких моделей - схемы параллельных программ, отражаюгце с разной степенью детализации структурные свойства программ; графовые модели и асинхронные сети (Петри сети), являюпиеся математич. абстракциями дискретных процессов и систем.



\begin{center}

\includegraphics[max width=\textwidth]{2023_07_08_5a2e49a87865251743a1g-1(2)}

\end{center}



\begin{center}

\includegraphics[max width=\textwidth]{2023_07_08_5a2e49a87865251743a1g-1}

\end{center}



\begin{center}

\includegraphics[max width=\textwidth]{2023_07_08_5a2e49a87865251743a1g-1(1)}

\end{center}



ПРОГРАММИРОВАНИЕ ТЕОРЕТИЧЕСКОЕ математическая дисцишлина, изучающая математич. абстракции программ, трактуемых как объекты, выраженные на формальном языке, обладающие определенной информационной и логич. структурой и подлежащие исполнению на автоматич. устройствах. П. т. сформировалось преимущественно на основе двух моделей вычислений: последовательных программ с памятью, или операторных программ, и рекурсивных программ. Обе модели строятся над нек-рой абстрактной алгебраич. системой $\langle D, \Phi, \Pi\rangle$, образованной предметной областью $D$, конечным набором (сигнатурой) функциональных $\Phi=\left\{\varphi_{1}, \ldots, \varphi_{m}\right\}$ и предикатных $\Pi=\left\{\pi_{1}, \ldots, \pi_{n}\right\}$ символов с заданным для каждого символа числом его аргументов (а р н о с т ь ю).



Определение класса программ слагается из трех частей: схемы программы (синтаксиса), интерпретации и семантики. С х е м п п о г р а м м ы- әто конструктивный объект, показывающий, как строится программа с использованием сигнатуры и других формальных символов. И н т е $\mathrm{p}$ п р е т а ц и я - это задание конкретной предметной области и сопоставление символам сигнатуры конкретных функций и предикатов (базовых операций), согласованных с предметной областью и арностью символов. С е м а н т и к аэто способ сопоставления каждой программе результата ее выполнения. Как правило, с программами связывают вычисляемые ими функции. Интерпретация обычно входит в семантику как параметр, поэтому схема программы задает множество программ и вычисляемых ими функций, к-рое получается при варьировании интерпретаций над нек-рым запасом базовых операций.



Схе ма п р о г р а м ы с п а м я т ью, называемая также ал голо по д о бно й, или о п ер а т о р в о й, с х е м о й, задается в виде конечного ориентированного г р а ф а п е р е х о д о в, имеющего обычно одну входную и одну выходную вершины, вершины с одной (преобразователи) и двумя (распознаватели) исходящими дугами. С помощью символов сигнатуры и счетного множества символов переменных и констант обычным образом строится множество функциональных и предикатных термов. Каждому распознавателю сопоставляется нек-рый предикатный терм, а преобразователю - оператор присваивания, имеющий вид $y:=\tau$, где $y$ - символ переменной, а $\tau$ - функциональный терм. Конечная совокупность всех переменных в схеме образует ее п а м я т ь. Интерпретация в дополнение к конкретизации базовых операций предписывает каждой переменной область ее изменения, а каждой константе - ее значение. Для программ с памятью наиболее обычна т. н. о п е р а ц и о н н а я с е м а н т и к а, состоящая из алгоритма выполнения программы на заданном состоянии памяти. Программа выполняется при движении по графу переходов. При попадании на распознаватель вычисляется предикатный терм и происходит переход по дуге, соответствующей значению предиката. При попадании на преобразователь с оператором $y:=\tau$ вычисляется значение $\tau$ и присваивается переменной $y$. Результат выполнения программы - состояние памяти при попадании на выходную вершину.



Схем а реку р с и вно й прог ра.ммы, или р е к у р с и в н а я с х е м а, использует кроме функциональных т. н. у с ло в ны е те р мы, образующие вместе с первыми множество вычислительных термов. Условный терм задает вычисление методом разбора случаев и имеет вид $\left(\pi\left|\tau_{1}\right| \tau_{2}\right)$, где $\pi-$ предикатный, а $\tau_{1}$ и $\tau_{2}$ - вычислительные термы, и соответствует конструкции условного выражения в алголе: if $\pi$ then $\tau_{1}$ else $\tau_{2}$.



Термы строятся над счетным множеством входных и формальных переменных, констант и символов опреде-





\end{document}